% =============================================================================
% GitHub Actions 兼容的学术论论文模板
% =============================================================================

\documentclass[12pt, a4paper]{article}

% =============================================================================
% 宏包加载
% =============================================================================

% 字体与编码
\usepackage[utf8]{inputenc}
\usepackage[T1]{fontenc}
\usepackage{mathptmx}
\usepackage{helvet}

% 语言支持
\usepackage[english]{babel}

% 页面布局
\usepackage[a4paper, margin=2.5cm]{geometry}
\usepackage{parskip}

% 图表支持
\usepackage{graphicx}
\usepackage{float}
\usepackage{caption}
\usepackage{subcaption}

% 表格增强
\usepackage{booktabs}
\usepackage{tabularx}
\usepackage{multirow}

% 数学公式
\usepackage{amsmath, amssymb, amsfonts}
\usepackage{mathtools}
\usepackage{physics}

% 单位
\usepackage{siunitx}

% 代码高亮
\usepackage{listings}
\usepackage{xcolor}

% 算法伪代码
\usepackage[ruled,vlined,linesnumbered]{algorithm2e}

% 参考文献设置
\usepackage[backend=biber,style=ieee]{biblatex}
\addbibresource{references.bib}

% 超链接
\usepackage{hyperref}
\hypersetup{
    colorlinks=true,
    linkcolor=blue,
    citecolor=blue,
    urlcolor=blue,
}

% 交叉引用
\usepackage{cleveref}
\crefname{section}{Section}{Sections}
\crefname{figure}{Figure}{Figures}
\crefname{table}{Table}{Tables}

% =============================================================================
% 代码高亮设置
% =============================================================================
\definecolor{codegreen}{rgb}{0,0.6,0}
\definecolor{codegray}{rgb}{0.5,0.5,0.5}
\definecolor{codepurple}{rgb}{0.58,0,0.82}
\definecolor{backcolour}{rgb}{0.95,0.95,0.92}

\lstdefinestyle{mystyle}{
    backgroundcolor=\color{backcolour},
    commentstyle=\color{codegreen},
    keywordstyle=\color{magenta},
    numberstyle=\tiny\color{codegray},
    stringstyle=\color{codepurple},
    basicstyle=\ttfamily\footnotesize,
    breakatwhitespace=false,
    breaklines=true,
    captionpos=b,
    keepspaces=true,
    numbers=left,
    numbersep=5pt,
    showspaces=false,
    showstringspaces=false,
    showtabs=false,
    tabsize=2
}
\lstset{style=mystyle}

% =============================================================================
% 标题信息
% =============================================================================
\title{Academic Paper Template\\
\large GitHub Actions Compatible}
\author{Author Name}
\date{\today}

% =============================================================================
% 文档开始
% =============================================================================
\begin{document}

\maketitle

\begin{abstract}
This is an academic paper template compatible with GitHub Actions. It compiles automatically when you push to the main branch and generates a PDF that is published to GitHub Releases.

Keywords: LaTeX, GitHub Actions, Continuous Integration
\end{abstract}

\newpage
\tableofcontents
\newpage

\section{Introduction}

This template demonstrates automatic LaTeX compilation using GitHub Actions~\cite{example2024}.

Features include:
\begin{itemize}
    \item Automatic compilation on push
    \item PDF generation and release
    \item Code highlighting
    \item Algorithm support
\end{itemize}

\section{Methods}

\subsection{Mathematical Formulas}

Einstein's mass-energy equation $E=mc^2$.

Schrödinger equation:
\begin{equation}
    i\hbar\frac{\partial}{\partial t}\Psi(\mathbf{r},t) = \hat{H}\Psi(\mathbf{r},t)
    \label{eq:schrodinger}
\end{equation}

\subsection{Algorithm}

\begin{algorithm}
    \SetAlgoLined
    \KwIn{Input parameters}
    \KwOut{Output results}
    Initialize variables\;
    \While{condition is met}{
        Execute operation\;
    }
    Return results\;
    \caption{Example Algorithm}
    \label{algo:example}
\end{algorithm}

\subsection{Code Example}

\begin{lstlisting}[language=Python, caption=Python Code Example]
def hello_world():
    print("Hello, World!")
    return 0
\end{lstlisting}

\section{Results}

\begin{table}[H]
    \centering
    \caption{Experimental Results}
    \label{tab:results}
    \begin{tabular}{lccc}
        \toprule
        Method & Accuracy & Precision & F1 Score \\
        \midrule
        Method A & 85.3 & 82.1 & 83.7 \\
        Method B & 87.6 & 84.5 & 86.0 \\
        \textbf{Our Method} & \textbf{90.2} & \textbf{88.7} & \textbf{89.4} \\
        \bottomrule
    \end{tabular}
\end{table}

Results from \cref{tab:results} show that our method performs best.

\section{Conclusion}

This paper demonstrates a GitHub Actions-compatible LaTeX template.

\newpage
\printbibliography

\end{document}
